%============================================================
% BAB 1: PENDAHULUAN
% Struktur:
%   Proposal/Progres (\ifbukupa = false):
%     1.1 Latar Belakang, 1.2 Permasalahan, 1.3 Tujuan,
%     1.4 Manfaat, 1.5 Sistematika Penulisan
%   Buku PA (\ifbukupa = true):
%     1.1 Latar Belakang, 1.2 Permasalahan, 1.3 Batasan Masalah,
%     1.4 Tujuan, 1.5 Manfaat, 1.6 Sistematika Penulisan
%============================================================
\chapter{PENDAHULUAN}

\section{LATAR BELAKANG}

Deskripsikan latar belakang dari permasalahan yang akan diangkat pada penelitian proyek
akhir. Latar belakang berisi penjelasan dari \textit{Problem Domain} yang termuat pada
judul penelitian. Misalkan penulis mengambil suatu judul ``Deteksi Penyakit Kanker
Dengan Sistem Pakar Berbasis Fuzzy'', maka \textit{Problem Domain}-nya adalah tentang
penyakit kanker, sehingga penulis dapat menceritakan tingkat urgensinya seperti
kenaikan jumlah penderita dari tahun ke tahun. Latar belakang yang baik berisikan
\textit{problem domain} yang mempunyai tingkat urgensi tinggi.

% Contoh kutipan tidak langsung (parafrase):
% Menurut penelitian sebelumnya, hal ini sangat penting dilakukan \cite{}.

\section{PERMASALAHAN}

Deskripsikan dengan jelas permasalahan yang ingin diteliti pada proyek akhir.
Permasalahan berisi penjelasan dari \textit{Problem} yang termuat pada judul penelitian.
Deskripsi masalah sebaiknya dituliskan dengan gaya bahasa deskriptif. Uraian
permasalahan yang baik manakala penulis berhasil meyakinkan pembaca tentang seberapa
tinggi tingkat urgensi dari permasalahan tersebut, sehingga pembaca yakin bahwa
permasalahan tersebut membutuhkan solusi dari penelitian proyek akhir.

% Batasan Masalah: hanya tampil di Buku PA
\ifbukupa
\section{BATASAN MASALAH}

Deskripsikan dengan jelas hal-hal yang perlu dibatasi dari permasalahan yang diselesaikan
dalam proyek akhir. Deskripsi batasan masalah sebaiknya dituliskan dalam bentuk
\textit{itemik}.

\begin{enumerate}[label=\arabic*., leftmargin=*]
  \item Batasan masalah pertama.
  \item Batasan masalah kedua.
  \item Batasan masalah ketiga.
\end{enumerate}
\fi

\section{TUJUAN}

Deskripsikan dengan jelas tujuan penelitian proyek akhir yang diangkat. Tujuan proyek
akhir harus jelas, singkat, dan mengandung klaim orisinalitas. Tuliskan secara
argumentatif apa saja fitur-fitur yang ditawarkan pada penelitian sebagai solusi baru
untuk mengklaim orisinalitas. Tujuan sebaiknya diawali dengan kalimat pembuka seperti:

``Penelitian proyek akhir ini mengajukan suatu pendekatan baru untuk \ldots dengan
mempresentasikan \ldots''. Kalimat-kalimat selanjutnya kemudian memperjelas
fitur-fitur unik yang ditawarkan pada penelitian sebagai suatu solusi untuk menjawab
permasalahan.

\section{MANFAAT}

Uraikan kontribusi proyek akhir pada pengembangan ilmu pengetahuan teknologi dan seni,
pemecahan masalah pembangunan, atau pengembangan kelembagaan. Kontribusi sebaiknya
bersifat spesifik, tidak terlalu luas, dan tidak terkesan mengada-ada. Jelaskan siapa
yang mendapatkan manfaat dari penelitian dan dalam bentuk apa manfaatnya.

\section{SISTEMATIKA PENULISAN}

Jelaskan tentang sistematika pembahasan dalam buku proyek akhir yang meliputi:

% File sistematika ditentukan oleh runner (main_*.tex) via \SistematikaFile.
% Proposal: pages/sistematika_3bab | Progres/Buku: pages/sistematika_5bab
\input{\SistematikaFile}
